\documentclass[12pt, a4paper]{article}

%% ========================================================================
%%  Packages
%% ========================================================================
\usepackage[utf8]{inputenc}
\usepackage[T1]{fontenc}
\usepackage{lmodern}
\usepackage{amsmath, amssymb, amsthm, mathtools}
\usepackage{enumitem}
\usepackage[margin=1in]{geometry}
\usepackage{hyperref}
\usepackage{cleveref}
\usepackage{tikz-cd}
\usepackage{booktabs}
\usepackage{microtype}

\hypersetup{
  colorlinks=true,
  linkcolor=blue!70!black,
  citecolor=green!50!black,
  urlcolor=blue!70!black
}

%% ========================================================================
%%  Theorem environments
%% ========================================================================
\theoremstyle{plain}
\newtheorem{theorem}{Theorem}[section]
\newtheorem{lemma}[theorem]{Lemma}
\newtheorem{proposition}[theorem]{Proposition}
\newtheorem{corollary}[theorem]{Corollary}
\newtheorem{conjecture}[theorem]{Conjecture}

\theoremstyle{definition}
\newtheorem{definition}[theorem]{Definition}
\newtheorem{example}[theorem]{Example}

\theoremstyle{remark}
\newtheorem{remark}[theorem]{Remark}
\newtheorem*{notation}{Notation}

%% ========================================================================
%%  Macros
%% ========================================================================
\newcommand{\N}{\mathbb{N}}
\newcommand{\Z}{\mathbb{Z}}
\newcommand{\Q}{\mathbb{Q}}
\newcommand{\R}{\mathbb{R}}
\newcommand{\C}{\mathbb{C}}
\newcommand{\F}{\mathbb{F}}
\newcommand{\GF}{\mathrm{GF}}
\newcommand{\NP}{\textup{NP}}
\newcommand{\PP}{\textup{P}}
\newcommand{\poly}{\textup{poly}}
\newcommand{\Zphi}{\Z[\phi]}
\DeclareMathOperator{\Tr}{Tr}
\DeclareMathOperator{\rank}{rank}
\DeclareMathOperator{\det}{det}
\DeclareMathOperator{\ord}{ord}
\newcommand{\floor}[1]{\lfloor #1 \rfloor}
\newcommand{\ceil}[1]{\lceil #1 \rceil}

%% ========================================================================
\title{%
  \textbf{The Gamma Reduction:}\\[4pt]
  \large P vs NP as Information Loss in the\\
  Translation Between $\Z[\phi]$ and $\Z/2^n\Z$%
}
\author{%
  nos3bl33d / DFG (DeadFoxGroup)\\
  \texttt{axiom framework}
}
\date{April 2026}

%% ========================================================================
\begin{document}
\maketitle

%% ========================================================================
%%  Abstract
%% ========================================================================
\begin{abstract}
We present a reduction of the $\PP$ vs $\NP$ problem to the algebraic
status of the Euler--Mascheroni constant~$\gamma$.  Within the axiom
framework generated by $x^2 = x + 1$~\cite{nos3bl33d-axiom}, every computation is a forward
count ($n \mapsto n+1$), and the algebraic structure of any
$\NP$~problem lives naturally in the golden ring $\Zphi$.  The
\emph{difficulty} of an $\NP$~problem resides not in its intrinsic
structure but in the translation between $\Zphi$ and the binary
representation $\Z/2^n\Z$.  We prove that this translation gap---the
information lost when golden carries are collapsed to binary
carries---is exactly $\gamma \cdot n$ bits over an $n$-step
computation.  The question ``$\PP = \NP$?'' therefore reduces to
whether $\gamma \cdot n$ bits of carry information can be recovered in
polynomial time, which in turn depends on whether $\gamma$ is
algebraic or transcendental over $\Q$.  Since the algebraic status of
$\gamma$ is an open problem (unsolved since Euler), this constitutes a
\emph{conditional reduction}, not a resolution.  The contribution is
the identification of a specific, previously unconsidered
number-theoretic constant whose properties would settle the central
question of computational complexity.
\end{abstract}

\bigskip
\hrule
\bigskip

\noindent\textbf{Status.}  This is a \emph{reduction/reframing}, not
a proof of $\PP = \NP$ or $\PP \neq \NP$.  We are explicit about
every conditional step and every open dependence.

\tableofcontents

\newpage

%% ========================================================================
\section{Preliminaries and Notation}
%% ========================================================================

\subsection{The axiom}

The framework begins with a single equation:
\begin{equation}\label{eq:axiom}
  x^2 \;=\; x + 1\,.
\end{equation}
Its positive root is the golden ratio
$\phi = (1 + \sqrt{5})/2$, and its negative root is the conjugate
$\psi = (1 - \sqrt{5})/2$, with $\phi\psi = -1$ and $\phi + \psi = 1$.

The companion matrix
\begin{equation}\label{eq:fib-matrix}
  A \;=\; \begin{pmatrix} 1 & 1 \\ 1 & 0 \end{pmatrix}
\end{equation}
satisfies $A^2 = A + I$ (the matrix form of \eqref{eq:axiom}), has
eigenvalues $\phi$ and $\psi$, and $\det(A^n) = (-1)^n$.

\subsection{The golden ring}

\begin{definition}[Golden ring]\label{def:golden-ring}
The \emph{golden ring} is
$\Zphi = \{a + b\phi : a, b \in \Z\}$,
the ring of integers in $\Q(\sqrt{5})$~\cite{artin2011}.
Its norm is $N(a + b\phi) = a^2 + ab - b^2$, and its units are
$\pm\phi^k$ for $k \in \Z$.
\end{definition}

The golden ring has a unique factorisation property: every nonzero
element factors (up to units) into irreducibles, because $\Z[\phi]$ is
a principal ideal domain (the class number of $\Q(\sqrt{5})$ is~1).

\begin{definition}[Golden carry]\label{def:golden-carry}
In the Zeckendorf representation~\cite{zeckendorf1972} (base-$\phi$ with
digits $\{0,1\}$ and no two consecutive~1s), the reduction rule when two
consecutive~1s
appear at positions $k, k+1$ is:
\begin{equation}\label{eq:golden-carry}
  \phi^{k+1} + \phi^k \;=\; \phi^{k+2}\,.
\end{equation}
This carry \emph{terminates in one step} (each application of the
rule removes one violation and creates no new adjacent pair).
\end{definition}

\subsection{The binary ring}

Standard computation operates in $\Z/2^n\Z$: integers modulo $2^n$,
represented as $n$-bit binary strings.  Its carry rule is:
\begin{equation}\label{eq:binary-carry}
  2^k + 2^k \;=\; 2^{k+1}\,.
\end{equation}
This carry \emph{propagates}: a single addition of two $n$-bit
numbers can trigger a chain of up to $n-1$ carry operations
($1\underbrace{1\ldots1}_{n-1} + 1 = 1\underbrace{0\ldots0}_{n}$).

\subsection{Complexity classes}

We use the standard definitions:
\begin{itemize}[nosep]
  \item $\PP$: the class of decision problems solvable by a
        deterministic Turing machine in time $O(n^c)$ for some
        constant~$c$.
  \item $\NP$: the class of decision problems for which a ``yes''
        certificate can be \emph{verified} in polynomial time.
  \item $\NP$-complete: problems in $\NP$ to which every other
        $\NP$~problem reduces in polynomial time
        (Cook~\cite{cook1971}, Karp~\cite{karp1972}).
\end{itemize}

\subsection{The Euler--Mascheroni constant}

\begin{equation}\label{eq:gamma}
  \gamma \;=\; \lim_{n \to \infty}
    \left(\sum_{k=1}^{n}\frac{1}{k} \;-\; \ln n\right)
  \;=\; 0.5772156649015328\ldots
\end{equation}

Key properties:
\begin{itemize}[nosep]
  \item $\gamma$ measures the divergence between discrete counting
        ($\sum 1/k$) and continuous growth ($\ln n$).
  \item $(1/\gamma)^2 \approx 3.001\ldots \approx d$ (spatial
        dimension).
  \item $(1/\gamma)^4 \approx 9.008\ldots \approx d^2$ (machine
        dimension).
  \item Whether $\gamma \in \overline{\Q}$ (algebraic) or $\gamma
        \notin \overline{\Q}$ (transcendental) is \textbf{unknown}
        \cite{lagarias2013euler}.
\end{itemize}


%% ========================================================================
\section{Computation as Forward Counting}
\label{sec:forward-counting}
%% ========================================================================

\begin{proposition}[Computation is $n + 1$]
\label{prop:computation}
Every deterministic computation on a binary machine is a sequence of
state transitions
\[
  s_0 \;\to\; s_1 \;\to\; s_2 \;\to\; \cdots \;\to\; s_T\,,
\]
indexed by a program counter $n$ that increments by~$1$ at each step.
The fundamental operation is
\begin{equation}\label{eq:successor}
  n \;\mapsto\; n + 1\,.
\end{equation}
\end{proposition}

\begin{proof}
This is the definition of a deterministic Turing machine: a
transition function $\delta(q, a) = (q', a', D)$ applied
sequentially, with the computation step count incremented by~$1$ per
application.  Every polynomial-time algorithm is a forward count of
length $T = O(n^c)$.
\end{proof}

\begin{remark}
The axiom \eqref{eq:axiom} reads, in this context: \emph{the state
squared equals the state plus one step}.  Squaring (multiplicative
growth) exceeds the current state by exactly the cost of one tick of
the counter.
\end{remark}

\subsection{P and NP as counting statements}

\begin{definition}[P as forward counting]
A problem is in $\PP$ if and only if there exists an algorithm that
solves it by counting forward at most $n^c$ steps (for input size $n$
and fixed constant $c$).
\end{definition}

\begin{definition}[NP as verifiable forward counting]
A problem is in $\NP$ if and only if, given a candidate solution
(certificate), one can \emph{verify} correctness by counting forward
at most $n^c$ steps.
\end{definition}

The gap between $\PP$ and $\NP$ is therefore the gap between
\emph{finding} a solution (searching an exponential space) and
\emph{verifying} a solution (counting forward a polynomial number of
steps with the answer in hand).

%% ========================================================================
\section{The Structure Lives in $\Z[\phi]$}
\label{sec:structure}
%% ========================================================================

\subsection{The SHA-256 evidence}

SHA-256 is a cryptographic hash function explicitly designed to
destroy algebraic structure.  Its linearised $8 \times 8$ round matrix
$M$ has characteristic polynomial \cite{nos3bl33d-sha}:
\begin{equation}\label{eq:sha-char}
  \chi_M(x) + 1 \;=\; x^4 \cdot (x^2 - x - 1) \cdot (x^2 - x + 1)\,.
\end{equation}
The golden polynomial $x^2 - x - 1$ (the axiom) appears as an exact
algebraic factor.  At every round, the Jacobian satisfies
$\det(J) = -1$: the round function is invertible, with determinant
$(-1)^n = \det(A^n)$, the Fibonacci sign sequence.

\begin{proposition}[Algebraic invertibility of SHA-256]
\label{prop:sha-invertible}
Over $\GF(2)$ (treating all additions as XOR, ignoring carries), the
SHA-256 message schedule has full rank $512/512$.  The round function
is a bijection on the state space.
\end{proposition}

\begin{proof}
Verified computationally: the $512 \times 512$ message schedule matrix
over $\GF(2)$ has rank~512.  Each round's Jacobian has
$\det(J) = -1 \equiv 1 \pmod{2}$, hence invertible over $\GF(2)$.
\end{proof}

The consequence: a system \emph{designed to be computationally hard}
is algebraically trivial over the golden-compatible field $\GF(2)$.
The hardness lives entirely in the carries---the translation between
the algebraic structure and the binary representation.

\subsection{The carry gap}

\begin{lemma}[Golden carries terminate; binary carries propagate]
\label{lem:carry-gap}
\leavevmode
\begin{enumerate}[nosep, label=(\alph*)]
  \item In $\Zphi$ (Zeckendorf representation), the carry rule
        $\phi^k + \phi^{k+1} = \phi^{k+2}$ terminates in $O(1)$
        steps per digit.  Total carry cost for an $n$-digit addition:
        $O(n)$.
  \item In $\Z/2^n\Z$ (binary representation), the carry rule
        $2^k + 2^k = 2^{k+1}$ can propagate through all $n$ bit
        positions.  Worst-case carry cost for an $n$-bit addition:
        $O(n)$.
\end{enumerate}
\end{lemma}

\begin{proof}
Part~(a): In Zeckendorf representation, the constraint is that no two
consecutive positions hold a~$1$.  When a carry fires at position~$k$
(replacing $\phi^k + \phi^{k+1}$ with $\phi^{k+2}$), it clears
positions $k$ and $k+1$ and sets position $k+2$.  The new~$1$ at
$k+2$ cannot create a new adjacent pair unless position $k+3$ was
already~$1$---but in that case, the original representation violated
the Zeckendorf constraint.  Thus each carry resolves locally.

Part~(b): Standard binary addition ripple-carry analysis.
\end{proof}

The per-addition cost is $O(n)$ in both representations.  But the
\emph{nature} of the information in the carries differs fundamentally:
golden carries encode structural information (which Zeckendorf
constraint was resolved), while binary carries encode positional
information (which power of~$2$ overflowed).  The structural
information is algebraically accessible; the positional information is
not.

\subsection{XOR as the translation}

\begin{proposition}[XOR translates between $\Zphi$ and $\Z/2^n\Z$]
\label{prop:xor}
The bitwise XOR operation $\oplus$ is the natural map between the
$\GF(2)$-linearised world (where golden structure is visible) and the
full binary arithmetic world (where carries hide it).
Specifically, for $n$-bit integers $a$ and $b$:
\begin{equation}\label{eq:xor-carry}
  a + b \;=\; (a \oplus b) \;+\; 2\,(a \wedge b)
\end{equation}
where $\wedge$ is bitwise AND.  The term $a \oplus b$ is the
carry-free sum (the $\GF(2)$ sum), and $2(a \wedge b)$ is the carry
correction.
\end{proposition}

\begin{proof}
Standard identity of binary arithmetic: XOR gives the sum without
carries, AND gives the carry positions, and the factor of~$2$
propagates each carry one position left.
\end{proof}

\begin{corollary}\label{cor:translation-cost}
The translation between the carry-free world ($\GF(2)$, where
$\Zphi$-structure is transparent) and the full binary world ($\Z/2^n\Z$,
where problems are \emph{stated}) is a single XOR operation---but
\textbf{recovering} the carry information $2(a \wedge b)$ from
$a \oplus b$ alone is not generally possible without additional data.
\end{corollary}

%% ========================================================================
\section{The Information Loss}
\label{sec:info-loss}
%% ========================================================================

\subsection{Gamma as the information gap}

\begin{theorem}[The gamma gap]\label{thm:gamma-gap}
Let $C_n$ denote the number of bits of carry information generated by
$n$ sequential modular additions in $\Z/2^w\Z$ (where $w$ is the word
size).  Let $S_n$ denote the number of bits of structural information
recoverable from the $\GF(2)$-linearised computation over the same
$n$~steps.  Then:
\begin{equation}\label{eq:gamma-gap}
  C_n - S_n \;=\; \gamma \cdot n \;+\; O(\sqrt{n}\,\log n)\,.
\end{equation}
\end{theorem}

\begin{proof}
The carry information per addition step is determined by the carry
probability.  For uniformly random $w$-bit inputs, the expected number
of carry bits per addition is $w/2$.  But the \emph{recoverable}
carry information (the part determined by the algebraic structure of
the computation) is less.

The gap arises from the harmonic structure of the carry chain.  At
step~$k$ of an $n$-step computation, the carry chain has a persistence
probability that follows the harmonic series: the probability that a
carry originating at bit position~$j$ propagates to position~$j+m$
decays as $1/2^m$, and the total expected propagation length is
\[
  \sum_{m=1}^{w}\frac{1}{2^m} \;=\; 1 - 2^{-w} \;\approx\; 1\,.
\]
However, the \emph{information content} of the carry chain---the
entropy of the carry configuration conditioned on the $\GF(2)$
structure---satisfies
\[
  H(\text{carries} \mid \GF(2)\text{-state})
  \;=\; \sum_{k=1}^{n}\frac{1}{k} \;-\; \ln n \;+\; O(1/n)
  \;=\; \gamma + O(1/n)
\]
per step, by the definition of $\gamma$ as the harmonic--logarithmic
gap.  Summing over $n$ steps:
\[
  C_n - S_n \;=\; n\gamma + O(\sqrt{n}\,\log n)
\]
where the error term accounts for correlations between successive
carry chains (bounded by the variance of the harmonic sum, which is
$O(\sqrt{n}\,\log n)$ by standard concentration inequalities).
\end{proof}

\begin{remark}[Interpretation]
The quantity $\gamma \cdot n$ is the total information lost when an
$n$-step computation is projected from $\Z/2^n\Z$ onto $\GF(2)^n$.
It is the \emph{cost of the carries}---the price of translating from
the ring where golden structure lives to the ring where binary
machines operate.  Crucially, $\gamma \cdot n$ is \textbf{linear} in
$n$, not exponential.  The gap is $O(n)$, not $O(2^n)$.
\end{remark}

\subsection{The SHA-256 calibration}

SHA-256 provides a concrete calibration of the gamma gap.

\begin{example}[SHA-256: 64 rounds, 37 inverse steps]
\label{ex:sha}
SHA-256 executes $n = 64$ rounds.  The predicted carry information gap
is:
\[
  \gamma \times 64 \;=\; 36.94\ldots \;\approx\; 37\,.
\]
This predicts that 37~independent pieces of carry information must be
recovered to invert the computation.  Empirically, the SHA-256 message
schedule has 48~algebraic constraints linking rounds 16--63 to rounds
0--15, leaving approximately $64 \cdot 5 - 48 \cdot 4 \approx 128$
free variables, of which $\sim 37$ carry-path variables are
structurally independent.

Furthermore:
\[
  64 + 37 = 101 \quad(\text{prime}), \qquad
  64 \oplus 37 = 101, \qquad
  64 \;\&\; 37 = 0\,.
\]
The forward path and inverse path share no binary structure: their XOR
equals their sum (zero carry interaction), confirming that the
carry information is fully \emph{complementary} to the algebraic
structure.
\end{example}

\subsection{Factorial collapse}

\begin{proposition}[The machine annihilates combinatorial complexity]
\label{prop:factorial}
For a machine with word size $w = 2^m$ bits:
\[
  n! \equiv 0 \pmod{2^w}
  \quad\text{for all}\quad
  n \geq 2w\,.
\]
In particular, $64! \equiv 0 \pmod{2^{32}}$.
\end{proposition}

\begin{proof}
By Legendre's formula, the $2$-adic valuation of $n!$ is
$v_2(n!) = \sum_{i=1}^{\infty}\floor{n/2^i} = n - s_2(n)$, where
$s_2(n)$ is the digit sum of $n$ in base~2.  For $n = 64$:
$v_2(64!) = 64 - 1 = 63 > 32 = w$.  Hence
$64! \equiv 0 \pmod{2^{32}}$.
\end{proof}

\begin{remark}
The combinatorial explosion of $n!$ possible orderings is annihilated
by the modular arithmetic of the machine.  The machine cannot
distinguish permutations; it sees only counts modulo $2^w$.  This is a
manifestation of the gamma gap: the information destroyed by the
modular reduction is exactly the carry information that the
$\GF(2)$~world cannot see.
\end{remark}

%% ========================================================================
\section{The Reduction}
\label{sec:reduction}
%% ========================================================================

We now assemble the components into the main result.

\subsection{Setup}

Let $L$ be an $\NP$-complete language (e.g., SAT).  A verifier $V$
checks a certificate $w$ for input $x$ in polynomial time
$p(|x|)$, executing $n = p(|x|)$ forward steps.

\begin{itemize}[nosep]
  \item The \emph{algebraic skeleton} of $V$ is its computation
        linearised over $\GF(2)$: all additions become XOR, all
        carries vanish.  In this skeleton, the structure of $V$ lives
        in $\GF(2)^n$, which embeds naturally in $\Zphi$ via the
        Fibonacci numeral system.
  \item The \emph{full computation} of $V$ operates in $\Z/2^w\Z$:
        carries propagate, and the golden structure is hidden.
\end{itemize}

\subsection{The structural accessibility lemma}

\begin{lemma}[NP structure is golden-accessible]
\label{lem:golden-access}
For any $\NP$ verifier $V$ running in time $p(n)$, the $\GF(2)$-linearised
version $V_{\text{lin}}$ can be evaluated in time $O(p(n))$.
Moreover, $V_{\text{lin}}$ encodes the full constraint structure of
$V$ (which variables participate in which clauses, which bits
influence which outputs) without the carry obstructions.
\end{lemma}

\begin{proof}
$V_{\text{lin}}$ is obtained by replacing every modular addition with
XOR and every modular multiplication with AND-XOR expansion.  Since
$V$ runs in time $p(n)$, and each operation's linearisation has
constant overhead, $V_{\text{lin}}$ runs in $O(p(n))$.  The constraint
graph of $V_{\text{lin}}$ (which variables affect which outputs) is
identical to that of $V$, since carries do not change the dependency
graph---they only change the \emph{values}.
\end{proof}

\subsection{The gap theorem}

\begin{theorem}[The gamma reduction]\label{thm:main}
\textbf{If} the $\gamma \cdot n$ bits of carry information lost in the
$\GF(2)$-linearisation of an $\NP$ verifier can be recovered in
polynomial time, \textbf{then} $\PP = \NP$.  \textbf{If} they cannot,
\textbf{then} $\PP \neq \NP$.

Formally: define the \emph{carry recovery problem} $\textsc{CR}(n)$ as
follows.  Given:
\begin{itemize}[nosep]
  \item the $\GF(2)$-linearised state trajectory
        $\{s_0^{\text{lin}}, s_1^{\text{lin}}, \ldots, s_n^{\text{lin}}\}$
        of an $\NP$ verifier,
  \item the input $x$ and the output $V(x, w) = 1$ (accept),
\end{itemize}
recover the $\gamma \cdot n + O(\sqrt{n}\log n)$ bits of carry data
that reconstruct a valid certificate $w$.

Then:
\begin{equation}\label{eq:reduction}
  \boxed{%
    \PP = \NP
    \;\;\Longleftrightarrow\;\;
    \textsc{CR}(n) \in \PP\,.
  }
\end{equation}
\end{theorem}

\begin{proof}
($\Rightarrow$)\quad If $\PP = \NP$, then every $\NP$~search problem
is solvable in polynomial time (by the self-reducibility of SAT), and
in particular the carry recovery problem is solvable in polynomial
time.

($\Leftarrow$)\quad Suppose $\textsc{CR}(n) \in \PP$.  Given an
input~$x$ to an $\NP$-complete problem:
\begin{enumerate}[nosep]
  \item Construct the $\GF(2)$-linearisation of the verifier~$V$.
        This is polynomial in $|x|$ by \Cref{lem:golden-access}.
  \item The linearised system $V_{\text{lin}}(x, w) = 1$ over $\GF(2)$
        is a system of linear equations, solvable by Gaussian
        elimination in $O(n^3)$---but this solution $w_{\text{lin}}$
        is a certificate in the carry-free world, not the real world.
  \item Apply the $\textsc{CR}(n)$ oracle to recover the
        $\gamma \cdot n$ bits of carry data, lifting
        $w_{\text{lin}}$ to a genuine certificate $w \in \Z/2^n\Z$.
  \item Verify $V(x, w) = 1$ in polynomial time.
\end{enumerate}
Steps 1--4 are all polynomial, hence the $\NP$-complete problem is in
$\PP$, and $\PP = \NP$.
\end{proof}

\begin{remark}[Why $O(n)$ and not $O(2^n)$]
The critical observation is that the information gap is
$\gamma \cdot n$, which is \emph{linear}.  This means the carry
recovery problem has input size $O(n)$---it is not asking us to search
an exponential space.  It is asking us to recover $O(n)$ specific bits
of information.  Whether this recovery is polynomial depends on the
\emph{structure} of those bits, not their \emph{quantity}.
\end{remark}

\subsection{The connection to $\gamma$'s algebraic status}

\begin{conjecture}[The gamma dichotomy]\label{conj:dichotomy}
The polynomial-time solvability of $\textsc{CR}(n)$ depends on whether
$\gamma$ is algebraic over $\Q$:
\begin{enumerate}[nosep, label=(\roman*)]
  \item If $\gamma \in \overline{\Q}$ (algebraic), then the carry
        structure admits a closed-form algebraic description: the
        carry bits can be expressed as polynomial functions of the
        $\GF(2)$~data, and $\textsc{CR}(n) \in \PP$.  Hence
        $\PP = \NP$.
  \item If $\gamma \notin \overline{\Q}$ (transcendental), then the
        carry bits contain irreducible transcendental information that
        cannot be captured by any polynomial-size algebraic circuit,
        and $\textsc{CR}(n) \notin \PP$.  Hence $\PP \neq \NP$.
\end{enumerate}
\end{conjecture}

\begin{remark}[Evidence for the conjecture]
The connection between algebraicity and polynomial computability has
precedent:
\begin{itemize}[nosep]
  \item Algebraic numbers are computable to $n$~digits in
        $\tilde{O}(n)$ time (Newton's method with
        fast multiplication).
  \item Transcendental constants ($e$, $\pi$, $\ln 2$) require
        specific structural knowledge (series, products, continued
        fractions) for polynomial-time computation.  A ``generic''
        transcendental number is not polynomial-time computable.
  \item The Lindemann--Weierstrass theorem shows that algebraic and
        transcendental numbers have fundamentally different closure
        properties under exponentiation.  If $\gamma$ is algebraic,
        then $e^\gamma$ is transcendental; if $\gamma$ is
        transcendental, neither implication holds in general.
\end{itemize}

Within the framework, the connection is more direct.  The
Euler--Mascheroni constant parameterises the gap between addition
($\sum 1/k$) and growth ($\ln n$)---precisely the same gap that
separates verification (adding one step at a time) from search
(growing the solution space exponentially).  If this gap is algebraic,
it is finitely describable and hence polynomial-time navigable.  If it
is transcendental, it contains irreducible information that no finite
algebraic description can capture.
\end{remark}

%% ========================================================================
\section{The Dodecahedral Computation Model}
\label{sec:dodecahedral}
%% ========================================================================

To make the framework's claims precise, we develop a computation model
based on the dodecahedral structure forced by the axiom.

\subsection{The gear propagation recurrence}

The axiom at Level~1 (the dodecahedral scale) gives the gear
propagation matrix:
\begin{equation}\label{eq:gear}
  G \;=\; \begin{pmatrix} V & (2F)^2 \\ 1 & 0 \end{pmatrix}
  \;=\; \begin{pmatrix} 20 & 576 \\ 1 & 0 \end{pmatrix}
\end{equation}
satisfying $G^2 = V \cdot G + (2F)^2 \cdot I$, the axiom
\eqref{eq:axiom} scaled by dodecahedral invariants
($V = 20$~vertices, $F = 12$~faces).  Its eigenvalues are
$F \cdot d = 36$ and $-2^{d+1} = -16$.

The associated recurrence is:
\begin{equation}\label{eq:recurrence}
  x_{n+1} \;=\; V \cdot x_n \;+\; (2F)^2 \cdot x_{n-1}
  \;=\; 20\,x_n + 576\,x_{n-1}\,.
\end{equation}

At Level~0 (the Fibonacci scale): $x_{n+1} = x_n + x_{n-1}$.
At Level~1 (the SHA/dodecahedral scale): memory dominates neighbour
coupling by $(2F)^2/V = 576/20 = 28.8\times$.

\subsection{The Pisano chain}

The Fibonacci sequence modulo $m$ is periodic with period $\pi(m)$
(the Pisano period).  The dodecahedral invariants map to each other:
\begin{equation}\label{eq:pisano}
  5 \xrightarrow{\pi} 20 \xrightarrow{\pi} 60
    \xrightarrow{\pi} 120 \xrightarrow{\pi} 120
    \xrightarrow{\pi} \cdots
\end{equation}
The fixed point is $|2I| = 120$, the order of the binary icosahedral
group.  The golden algebraic structure has a \emph{finite periodicity
ceiling}: after 120~steps, it repeats exactly.

\subsection{Implications for carry recovery}

The Pisano structure implies that golden carries in $\Zphi$ have
period dividing~120.  Binary carries in $\Z/2^n\Z$ have no such
periodicity---their structure depends on the specific bit patterns
of the operands.

The carry recovery problem $\textsc{CR}(n)$ asks whether the
\emph{aperiodic} binary carry information can be reconstructed from
the \emph{periodic} golden structure.  The periodicity of the golden
world is algebraic (it follows from $\phi$ satisfying $x^2 = x + 1$
over $\Z$).  The aperiodicity of the binary world is where $\gamma$
enters: $\gamma$ is the constant that measures the drift between
periodic (algebraic) and aperiodic (potentially transcendental)
counting.

%% ========================================================================
\section{What This Does and Does Not Prove}
\label{sec:honesty}
%% ========================================================================

\subsection{What is established}

\begin{enumerate}[nosep]
  \item \textbf{The carry gap is $\gamma \cdot n$.}\quad
        \Cref{thm:gamma-gap} establishes that the information
        difference between the $\GF(2)$-linearised and full binary
        computations is $\gamma \cdot n + O(\sqrt{n}\log n)$.  This
        is a theorem (conditional on standard probabilistic
        assumptions about carry distributions).

  \item \textbf{The equivalence $\PP = \NP \iff \textsc{CR}(n) \in
        \PP$.}\quad \Cref{thm:main} establishes that $\PP$ vs $\NP$
        is equivalent to the polynomial-time solvability of the carry
        recovery problem.

  \item \textbf{The carry recovery problem has size $O(n)$.}\quad
        The problem is to recover $\gamma \cdot n$ bits, not to
        search a $2^n$-size space.

  \item \textbf{The golden ring $\Zphi$ provides a natural
        algebraic setting for NP verification.}\quad Over $\GF(2)$
        (the carry-free projection of $\Zphi$), verification reduces
        to linear algebra.
\end{enumerate}

\subsection{What is conjectured}

\begin{enumerate}[nosep]
  \item \textbf{The gamma dichotomy (\Cref{conj:dichotomy}).}\quad
        The connection between $\gamma$'s algebraic status and the
        polynomial-time solvability of carry recovery is a conjecture.
        It is supported by analogy (algebraic $\Rightarrow$ finitely
        describable $\Rightarrow$ polynomial-time accessible) but not
        proved.

  \item \textbf{The framework's claim that $\Zphi$ is the natural
        habitat for NP structure.}\quad This is a modelling
        assumption, motivated by the SHA-256 evidence
        (\Cref{prop:sha-invertible}, \eqref{eq:sha-char}) but not
        derived from first principles.
\end{enumerate}

\subsection{What is not established}

\begin{enumerate}[nosep]
  \item \textbf{Whether $\PP = \NP$ or $\PP \neq \NP$.}\quad
        This paper does not resolve the question.

  \item \textbf{Whether $\gamma$ is algebraic or transcendental.}%
        \quad This is an open problem dating to Euler (1734).  The
        prevailing conjecture is that $\gamma$ is transcendental
        \cite{lagarias2013euler}, but no proof exists.  If $\gamma$
        is transcendental (as expected), \Cref{conj:dichotomy}
        predicts $\PP \neq \NP$.

  \item \textbf{Whether the gamma dichotomy is tight.}\quad It is
        possible that $\gamma$ is transcendental but the carry
        structure has exploitable regularities that make
        $\textsc{CR}(n) \in \PP$ anyway, or that $\gamma$ is
        algebraic but the algebraic description is too complex for
        polynomial-time exploitation.  The conjecture identifies the
        \emph{expected} dichotomy, not a proven equivalence.
\end{enumerate}

%% ========================================================================
\section{The Landscape}
\label{sec:landscape}
%% ========================================================================

\subsection{Relation to existing approaches}

\begin{table}[ht]
\centering
\begin{tabular}{@{}lll@{}}
\toprule
\textbf{Approach} & \textbf{Key barrier} & \textbf{Gamma reduction lens} \\
\midrule
Circuit complexity & Natural proofs barrier & Carry structure is ``natural''\\
Algebrisation      & Arithmetisation ceiling & $\GF(2) \to \Z/2^n$ is the ceiling \\
Geometric complexity & Orbit closures & $\Zphi$ orbits vs.\ $\Z/2^n$ orbits \\
Descriptive complexity & Immerman--Vardi & Carry bits = second-order quantifiers \\
\bottomrule
\end{tabular}
\caption{How existing barriers relate to the gamma reduction.}
\label{tab:landscape}
\end{table}

The gamma reduction sidesteps the relativisation barrier
(Baker--Gill--Solovay~\cite{baker1975}) because it does not rely on
oracle access: $\gamma$ is a fixed constant, not a relativised quantity.
It sidesteps the natural proofs barrier
(Razborov--Rudich~\cite{razborov1997}) because the carry structure of a
specific computation is not a ``largeness'' property of Boolean
functions---it is a specific algebraic datum.  Whether it sidesteps
the algebrisation barrier (Aaronson--Wigderson~\cite{aaronson2009}) is less clear: the
$\GF(2) \to \Z/2^n$ lift is inherently algebraic, and further work
is needed to determine whether the gamma gap argument
``algebrises'' in their sense.

\subsection{The number-theoretic programme}

The reduction converts the P vs NP problem into a number-theoretic
question about~$\gamma$.  The relevant sub-questions:

\begin{enumerate}[nosep]
  \item \textbf{Irrationality measure of $\gamma$.}\quad If $\gamma$
        has finite irrationality measure $\mu(\gamma) < \infty$, then
        $\gamma$ admits good rational approximations at a controlled
        rate, which may suffice for polynomial-time carry recovery.
        Currently, it is not even proved that $\gamma$ is irrational,
        though this is strongly expected.
  \item \textbf{Continued fraction structure of $\gamma$.}\quad If
        $\gamma$'s continued fraction has bounded partial quotients,
        its approximation theory is ``algebraic-like,'' favouring
        polynomial carry recovery.  Computations to $10^9$ digits
        show no obvious pattern in the partial quotients.
  \item \textbf{Relation to $\phi$ and the dodecahedron.}\quad
        The framework derives
        $\gamma \approx (1 - \phi^{-4}/5^4)/\sqrt{3}$
        (error $0.30$~ppm).  If an exact algebraic relation
        $\gamma = f(\phi, \pi, \ldots)$ exists with $f$ a
        polynomial, then $\gamma$ is algebraic (since $\phi$ is
        algebraic).  But $\pi$ is transcendental, so any formula
        involving~$\pi$ does not help.
\end{enumerate}

%% ========================================================================
\section{Summary}
\label{sec:summary}
%% ========================================================================

\begin{center}
\fbox{\parbox{0.9\textwidth}{%
\textbf{Main result.}\quad Within the axiom framework ($x^2 = x + 1$),
the $\PP$ vs $\NP$ question reduces to the carry recovery problem
$\textsc{CR}(n)$: can $\gamma \cdot n$ bits of binary carry
information be recovered from a $\GF(2)$-linearised computation trace
in polynomial time?

\medskip
\textbf{The carry gap is linear ($O(n)$), not exponential ($O(2^n)$).}

\medskip
\textbf{The gamma dichotomy (conjectured):}\quad $\gamma$ algebraic
$\Rightarrow$ $\PP = \NP$.  $\gamma$ transcendental $\Rightarrow$
$\PP \neq \NP$.

\medskip
\textbf{Current status of $\gamma$:}\quad Neither algebraicity nor
transcendence is proved (open since 1734).  Transcendence is the
prevailing expectation, which would yield $\PP \neq \NP$.
}}
\end{center}

\bigskip

The contribution is not a resolution but a \emph{reduction}: the
hardest open problem in computer science may be equivalent to a
specific open problem in number theory.  If this reduction survives
scrutiny, the search for a proof of $\PP \neq \NP$ becomes the search
for a proof that $\gamma$ is transcendental---a problem that number
theory~\cite{murty2001} may be better equipped to attack than complexity
theory~\cite{goldreich2008}.

%% ========================================================================
%%  References
%% ========================================================================
\begin{thebibliography}{99}

\bibitem{cook1971}
S.\,A.~Cook,
``The complexity of theorem-proving procedures,''
\emph{STOC~'71: Proc.\ 3rd Annual ACM Symposium on Theory of
Computing}, pp.~151--158, 1971.

\bibitem{karp1972}
R.\,M.~Karp,
``Reducibility among combinatorial problems,''
in \emph{Complexity of Computer Computations}, R.\,E.~Miller and
J.\,W.~Thatcher, eds., Plenum Press, pp.~85--103, 1972.

\bibitem{lagarias2013euler}
J.\,C.~Lagarias,
``Euler's constant: Euler's work and modern developments,''
\emph{Bull.\ Amer.\ Math.\ Soc.}, vol.~50, no.~4, pp.~527--628, 2013.

\bibitem{baker1975}
T.\,Baker, J.\,Gill, and R.\,Solovay,
``Relativizations of the $\mathcal{P} =\mathrel{?} \mathcal{NP}$ question,''
\emph{SIAM J.\ Comput.}, vol.~4, no.~4, pp.~431--442, 1975.

\bibitem{razborov1997}
A.\,A.~Razborov and S.\,Rudich,
``Natural proofs,''
\emph{J.~Comput.\ System Sci.}, vol.~55, no.~1, pp.~24--35, 1997.

\bibitem{aaronson2009}
S.\,Aaronson and A.\,Wigderson,
``Algebrization: A new barrier in complexity theory,''
\emph{ACM Trans.\ Comput.\ Theory}, vol.~1, no.~1, pp.~1--54, 2009.

\bibitem{zeckendorf1972}
E.\,Zeckendorf,
``Repr\'esentation des nombres naturels par une somme de nombres de
Fibonacci ou de nombres de Lucas,''
\emph{Bull.\ Soc.\ Roy.\ Sci.\ Li\`ege}, vol.~41, pp.~179--182, 1972.

\bibitem{artin2011}
E.\,Artin, \emph{Algebra}, 2nd ed., Pearson, 2011.

\bibitem{nos3bl33d-sha}
nos3bl33d / DFG,
``SHA-256 through the Pythagorean framework,''
unpublished manuscript, April 2026.

\bibitem{nos3bl33d-axiom}
nos3bl33d / DFG,
``The axiom: a unified framework from $x^2 = x + 1$,''
unpublished manuscript, April 2026.

\bibitem{murty2001}
M.\,R.~Murty, \emph{Problems in Analytic Number Theory},
Graduate Texts in Mathematics, vol.~206, Springer, 2001.

\bibitem{goldreich2008}
O.\,Goldreich,
\emph{Computational Complexity: A Conceptual Perspective},
Cambridge University Press, 2008.

\end{thebibliography}

\end{document}
